\section{Model}\label{sec:model}

\para{Transducer.} \ultra is a streaming RNN-T~\cite{graves2012rnnt} with a Conformer
encoder~\cite{gulati2020conformer}, a stateless predictor~\cite{ghodsi2020stateless} and an
auxiliary CTC head. The front end computes 80-band log-mel features at 100\,Hz; a causal
two-layer convolutional stem subsamples by four to 25\,Hz at width $d=384$. Each Conformer block
is macaron-style (half-step FFN, attention, causal depthwise convolution with kernel 15, FFN),
with six heads of width 64, an FFN expansion of four, rotary position
embeddings~\cite{su2021roformer}, and pre-norm residuals closed by a single final LayerNorm (a
post-norm variant of our first model collapsed to near time-invariant output). The predictor
embeds the last two non-blank tokens (256 dimensions each) and projects them to a 320-wide
joiner over 4,096 SentencePiece units plus blank. The CTC head (1,024 units) is weighted \CtcWeight{} in
training and unused at inference.

\para{Streaming.} The encoder consumes 160\,ms chunks (four 25\,Hz frames). Attention in chunk
$c$ sees the current chunk and a bounded left context of 128 frames (\LeftCtxSec\,s) and nothing after
the chunk; the convolution is causal. Hence every encoder frame depends only on audio up to the
end of its chunk, and the streaming model is the training model, not an approximation of it: a
unit test perturbs all future log-mel frames and checks that past encoder frames do not change.
Per stream, the state is a ring of keys and values per block, the convolution tail and two
predictor tokens.

\para{Multirate encoder (S\_MR).} Dense Conformers spend most of their arithmetic in the middle of
the stack at the full 25\,Hz rate. The multirate variant keeps the same width and block type but
changes where the frames are: four lower blocks run at 25\,Hz; a learned projection maps each
pair of frames to one \MidRateHz\,Hz frame; ten middle blocks run at \MidRateHz\,Hz with half the
chunk and half the left context; a learned restoration maps each frame back to two 25\,Hz
positions, added to a skip from the lower blocks; and two final blocks run at 25\,Hz
(\cref{fig:arch}). Per 160\,ms chunk the dense encoder applies $16\times4=64$ block-frame
operations and the multirate encoder $6\times4+10\times2=44$, before the two rate-change
projections. Weight bytes are the same, so the saving is in arithmetic, not memory traffic; how
much of it survives depends on whether the engine is compute-bound, which on our CPUs it is
(\cref{sec:engine}). The joiner and CTC head always read 25\,Hz states,
so decoding is unchanged. An odd final frame in a chunk is paired with zeros, both in the
engine and, after a fix we describe in \cref{sec:training}, in batched training.

\begin{figure}[t]
\centering
\resizebox{\columnwidth}{!}{%
\begin{tikzpicture}[
  font=\footnotesize, >={Stealth[length=4pt]},
  blk/.style={draw, rounded corners=2pt, minimum height=6mm, inner sep=2pt, align=center},
  fast/.style={blk, fill=blue!10, minimum width=17mm},
  slow/.style={blk, fill=orange!18, minimum width=17mm},
  io/.style={blk, fill=gray!10, minimum width=17mm},
  op/.style={blk, fill=white, minimum width=13mm, dashed}]
\node[io] (mel) {log-mel\\100\,Hz};
\node[io, right=4mm of mel] (stem) {causal stem\\$\downarrow$4};
\node[fast, right=4mm of stem] (low) {4 blocks\\25\,Hz};
\node[op, right=4mm of low] (pair) {pair\\$\downarrow$2};
\node[slow, below=6mm of pair, xshift=-10mm] (mid) {10 blocks\\\MidRateHz\,Hz};
\node[op, left=4mm of mid] (rest) {restore\\$\uparrow$2};
\node[fast, left=4mm of rest] (fin) {2 blocks\\25\,Hz};
\node[io, left=4mm of fin] (dec) {joiner\\+ CTC};
\draw[->] (mel) -- (stem); \draw[->] (stem) -- (low); \draw[->] (low) -- (pair);
\draw[->] (pair) |- (mid); \draw[->] (mid) -- (rest); \draw[->] (rest) -- (fin);
\draw[->] (fin) -- (dec);
\draw[->, densely dotted] (low.south) -- ++(0,-3mm) -| node[near start, below, inner sep=1pt]{skip} (rest.north);
\end{tikzpicture}}
\caption{The multirate encoder (S\_MR). Ten of sixteen Conformer blocks run at half rate inside
each 160\,ms chunk; the dense model (S) runs all sixteen at 25\,Hz.}
\label{fig:arch}
\end{figure}
